\section*{Limitations}
\label{sec:limitations}
\paragraph{The Hard split was chosen using model failures.} Its tasks were selected because five model families failed them, so those families' Hard scores are biased downward. Systems outside the selection (Opus~5.5, Fable~5.1, Astra and the partner runs of Sonnet~5, Terra, Sol and Haiku~4.5) show the same spread, from 25\% to 96\%. With \nHard{} tasks, single-task differences are noise.

\paragraph{One attempt per task.} Each system ran once, and every run had at least one flaky test. The partner runs carry a single grade and were not re-graded by us.

\paragraph{Training data.} The public repositories and fixes predate the evaluated models. The protocol stops an agent from fetching a solution, not from having seen it in training. The never-public set reduces this concern but does not remove it.

\paragraph{What re-grading cannot catch.} The clean container rejects changes made outside the patch, not code inside it. The imitation SDKs on the private set passed both grades. Our audits match commands and files, so they give lower bounds. Instructions still name many of the symbols the tests use, a property of the original benchmark design that we measured but did not remove.

\section*{Ethics Statement}
The public tasks come from open-source repositories under their original licences. The private tasks come from organisations that consented to their use for evaluation, and their code is not released. The paper describes how agents faked dependencies and how a sandbox grader can be fooled, because benchmark operators need this to defend against it; the clean-container re-grade blocks both. No personal data about the engineers is recorded.
